% ============================================================
% MÓDULO 0: PROBABILIDAD Y PENSAMIENTO ESTOCÁSTICO PARA RL
% ============================================================
\section{M0 · Probabilidad}
\module{Módulo 0: Probabilidad}{El lenguaje de la incertidumbre: de las frecuencias a la esperanza matemática}

\begin{frame}{¿Por qué probabilidad antes de Reinforcement Learning?}
  \begin{columns}[T,onlytextwidth]
    \column{0.48\textwidth}
      \begin{block}{El problema en el mundo real}
        \begin{itemize}
          \item Los robots se resbalan en el suelo húmedo.
          \item Los sensores devuelven mediciones con ruido.
          \item Los tiempos de entrega nunca son exactamente iguales.
        \end{itemize}
      \end{block}
      \vspace{0.2cm}
      \small En el almacén de \textbf{OptiBot}, una misma acción no siempre produce el mismo resultado.
    \column{0.48\textwidth}
      \begin{alertblock}{La solución: Pensamiento Estocástico}
        \begin{itemize}
          \item No podemos predecir el futuro exacto.
          \item \textbf{Sí} podemos modelar la distribución de futuros posibles.
          \item Toda función de valor en RL es una \textbf{esperanza matemática}.
        \end{itemize}
      \end{alertblock}
  \end{columns}
  \vspace{0.4cm}
  \keyidea{En RL no buscamos certezas absolutas, sino maximizar el \textbf{valor esperado} de nuestras ganancias a largo plazo.}
\end{frame}

\begin{frame}{Variables Aleatorias y Distribuciones}
  \begin{columns}[c,onlytextwidth]
    \column{0.50\textwidth}
      \begin{block}{Definición intuitiva}
        Una \textbf{variable aleatoria discreta} $X$ es una regla que asigna un número a cada posible resultado de un experimento.
      \end{block}
      \vspace{0.25cm}
      \textbf{Propiedades obligatorias:}
      \begin{itemize}
        \item $0 \le P(X = x) \le 1$ para todo valor $x$.
        \item $\sum_x P(X = x) = 1$ (la certeza total).
      \end{itemize}
    \column{0.46\textwidth}
      \centering
      \textbf{Caso OptiBot: Tiempo de Entrega $T$}\\
      \vspace{0.2cm}
      \begin{tabular}{@{}ccc@{}}
        \toprule
        \textbf{Días ($t$)} & \textbf{$P(T = t)$} & \textbf{Histograma} \\
        \midrule
        1 día & $0.30$ & \tikz{\fill[Teal] (0,0) rectangle (1.2,0.25);} \\
        2 días & $0.50$ & \tikz{\fill[TealDark] (0,0) rectangle (2.0,0.25);} \\
        3 días & $0.20$ & \tikz{\fill[Coral] (0,0) rectangle (0.8,0.25);} \\
        \bottomrule
      \end{tabular}\\
      \vspace{0.2cm}
      \small Suma: $0.30 + 0.50 + 0.20 = 1.00$
  \end{columns}
\end{frame}

\begin{frame}{Probabilidad Condicional y Regla Total}
  \begin{columns}[T,onlytextwidth]
    \column{0.48\textwidth}
      \begin{block}{Probabilidad Condicional}
        ¿Cómo cambia la probabilidad de $B$ sabiendo que ocurrió $A$?
        \[
          P(B \mid A) = \frac{P(A \cap B)}{P(A)}
        \]
      \end{block}
      \small Si OptiBot sabe que el piso está mojado ($A$), la probabilidad de derrape ($B$) aumenta.
    \column{0.48\textwidth}
      \begin{alertblock}{Ley de Probabilidad Total}
        Si el espacio se divide en eventos disjuntos $A_1, A_2, \dots, A_n$:
        \[
          P(B) = \sum_{i=1}^n P(B \mid A_i) \cdot P(A_i)
        \]
      \end{alertblock}
      \small ¡Esta suma ponderada es el corazón exacto de las ecuaciones de Bellman!
  \end{columns}
  \vspace{0.4cm}
  \keyidea{La probabilidad condicional formaliza cómo la información actual condiciona las probabilidades del futuro.}
\end{frame}

\begin{frame}{El Concepto Rey del RL: Valor Esperado}
  \begin{columns}[c,onlytextwidth]
    \column{0.50\textwidth}
      \begin{beamercolorbox}[rounded=true,sep=0.8em]{block body}
        \centering
        {\Large $\mathbb{E}[X] = \sum_{x} x \cdot P(X = x)$}
      \end{beamercolorbox}
      \vspace{0.3cm}
      \begin{itemize}
        \item Es el promedio ponderado a largo plazo.
        \item \textbf{Cálculo manual en OptiBot:}\\
        Ganancia neta por paquete según tiempo $T$:
        \begin{itemize}
          \item 1 día ($+\$10$): prob $0.30$
          \item 2 días ($+\$5$): prob $0.50$
          \item 3 días ($-\$2$): prob $0.20$
        \end{itemize}
      \end{itemize}
    \column{0.46\textwidth}
      \begin{exampleblock}{Cálculo manual con Lápiz y Papel}
        \begin{align*}
          \mathbb{E}[\text{Ganancia}] &= (10)(0.30) + (5)(0.50) + (-2)(0.20) \\
          &= 3.0 + 2.5 - 0.4 \\
          &= \mathbf{\$5.10 \text{ por paquete}}
        \end{align*}
      \end{exampleblock}
      \small Aunque ningún paquete dé exactamente \$5.10, en 10,000 entregas ganaremos en promedio \$5.10 por paquete.
  \end{columns}
\end{frame}

\begin{frame}[fragile]{Verificación en Python: Ley de Grandes Números}
  \begin{columns}[T,onlytextwidth]
    \column{0.62\textwidth}
\begin{lstlisting}[style=rlcode]
import numpy as np

# Valores y distribucion teorica
ganancias = np.array([10.0, 5.0, -2.0])
probabilidades = np.array([0.30, 0.50, 0.20])

# Valor esperado teorico
E_teorico = np.sum(ganancias * probabilidades)

# Simulacion de 10,000 entregas de OptiBot
np.random.seed(42)
muestras = np.random.choice(ganancias, size=10000, p=probabilidades)
E_empirico = np.mean(muestras)

print(f"Teorico:  ${E_teorico:.2f}")
print(f"Empirico: ${E_empirico:.2f}")
# Salida: Teorico: $5.10 | Empirico: $5.08
\end{lstlisting}
    \column{0.35\textwidth}
      \begin{block}{Lección del código}
        \begin{itemize}
          \item Cuando no conocemos la distribución teórica, podemos \textbf{estimar la esperanza promediando muestras}.
          \item Este principio simple es la base de los métodos \textbf{Monte Carlo} y del aprendizaje por interacción.
        \end{itemize}
      \end{block}
  \end{columns}
\end{frame}
